\clearpage
\newcount\SavedTranslationPage
\SavedTranslationPage=\value{page}
\hypersetup{pageanchor=false}
\begingroup
\pagestyle{empty}\thispagestyle{empty}
\pdfbookmark[0]{Translation sources and terminology}{translation-sources-I}
\noindent{\Large\bfseries Translation sources and terminology\par}\medskip
\small\raggedright
Historical text: Carl Friedrich Gauss, \emph{Werke}, Volume I, Göttingen, 1863. Historical editorial contributions retain their own authorship and signatures. Originally German passages are rendered in modern German in the German reader; their historical originals are supplied separately.

These English and German editions are AI-translated using OpenAI GPT-6.1 Sol, with AI-produced typesetting. The following named human translations and writings are scholarly references for the terminology and usage specified below. Gauss's original-language text remains the mathematical source.

\par\medskip\noindent\textbf{H. Maser.}\par\nobreak\smallskip\noindent German translator and editor, \emph{Carl Friedrich Gauss' Untersuchungen über höhere Arithmetik}, Berlin: Julius Springer, 1889. Corresponding arithmetical passages and German mathematical terminology.\par
\url{https://archive.org/details/carlfriedrichga00gausgoog}
\par\medskip\noindent\textbf{James Caddall Morehead and Adam Miller Hiltebeitel.}\par\nobreak\smallskip\noindent Translators of Gauss's \emph{General Investigations of Curved Surfaces of 1827 and 1825}, Princeton University Library, 1902; Project Gutenberg 36856, TeX edition. Related English geometrical terminology and usage.\par
\url{https://www.gutenberg.org/ebooks/36856}
\par\medskip\noindent\textbf{Paul Taylor; Bernard Leak.}\par\nobreak\smallskip\noindent Taylor's English translation of Gauss's second proof of the fundamental theorem of algebra, with corrections by Leak, 1983. Related English algebraic terminology.\par
\url{https://www.paultaylor.eu/misc/gauss.html}
\par\medskip\noindent\textbf{Ulrich Görtz.}\par\nobreak\smallskip\noindent \emph{Algebra}, WS 2021/22, 3.1.1–3.1.3 and opening of 3.1.4; \emph{Lineare Algebra II}, SS 2021, 18.4. Contemporary German arithmetic and algebraic usage.\par
\url{https://math.ug/a-ws2122/sect0006.html}\par
\url{https://math.ug/la2-ss21/sec-quotient-ring.html}
\par\medskip\noindent\textbf{Nikolay Moshchevitin.}\par\nobreak\smallskip\noindent \emph{Eine Bemerkung über positiv definite quadratische Formen und rationale Punkte}, arXiv:1510.06214v3. Related contemporary German terminology for quadratic forms.\par
\url{https://arxiv.org/abs/1510.06214v3}


\noindent\textbf{Publication and rights.} These editions are available free of charge for non-commercial, purely educational purposes. Publication rests on the publisher's presumption that the respective rights holders would not object; this is not established permission. Free access to a source does not itself authorize derivative publication. Rights in the originals and any protected derivative contributions remain with their respective rights holders. No project-wide licence or public-domain dedication is granted here. Rights holders may raise objections through the repository's issue page. The publisher undertakes to take down the affected material upon a rights holder's request.\par
\url{https://github.com/KokunoYumeto/modern-latex-manuscripts/issues}

\clearpage\endgroup\clearpage
\setcounter{page}{\the\SavedTranslationPage}
\hypersetup{pageanchor=true}
